% Generated by mm41 --commands-tex from the dispatch table; do not edit by hand.
\subsection*{Module \texttt{core} 1.3 --- arithmetic, functions, logic, the stack, words and the session}
\subsubsection*{Arithmetic}
\begin{longtable}{@{}>{\raggedright\arraybackslash}p{0.12\textwidth}>{\raggedright\arraybackslash}p{0.19\textwidth}>{\raggedright\arraybackslash}p{0.36\textwidth}>{\raggedright\arraybackslash}p{0.23\textwidth}@{}}
\toprule \textbf{Command} & \textbf{Stack} & \textbf{Description} & \textbf{Example}\\ \midrule \endfirsthead
\toprule \textbf{Command} & \textbf{Stack} & \textbf{Description} & \textbf{Example}\\ \midrule \endhead
\bottomrule \endlastfoot
\texttt{+} & \texttt{y x -{}-{} y+x} & add scalars or matrices; a scalar added to a matrix is added elementwise & \texttt{3 4 +}\\
\texttt{-{}} & \texttt{y x -{}-{} y-{}x} & subtract scalars or matrices (a scalar elementwise) & \texttt{10 4 -{}}\\
\texttt{*} & \texttt{y x -{}-{} y*x} & multiply scalars, or matrix product for matrices & \texttt{[2 2 \$ 1 2 3 4] [2 1 \$ 1 1] *}\\
\texttt{/} & \texttt{y x -{}-{} y/x} & divide scalars, or a matrix by a scalar; division by zero is an error & \texttt{1 3 /}\\
\texttt{\^{}} & \texttt{y x -{}-{} y\^{}x} & power; a square matrix to an integer power is a matrix power; a negative base with a fractional exponent promotes to complex & \texttt{2 10 \^{}}\\
\texttt{pow} & \texttt{y x -{}-{} y\^{}x} & power (alias of \^{}) & \texttt{2 10 pow}\\
\texttt{.*} & \texttt{A B -{}-{} A.*B} & elementwise matrix product & \texttt{[2 2 \$ 1 2 3 4] [2 2 \$ 1 1 2 2] .*}\\
\texttt{./} & \texttt{A B -{}-{} A./B} & elementwise division (matrix or scalar operands) & \texttt{[2 2 \$ 1 2 3 4] [2 2 \$ 1 2 1 2] ./}\\
\texttt{.\^{}} & \texttt{A B -{}-{} A.\^{}B} & elementwise power (matrix or scalar operands) & \texttt{[2 2 \$ 1 2 3 4] 2 .\^{}}\\
\texttt{chs} & \texttt{x -{}-{} -{}x} & change sign (elementwise for matrices) & \texttt{5 chs}\\
\texttt{inv} & \texttt{x -{}-{} 1/x} & reciprocal (elementwise for matrices; see minv for the matrix inverse) & \texttt{4 inv}\\
\texttt{abs} & \texttt{z -{}-{} \textbar{}z\textbar{}} & absolute value; magnitude for complex; elementwise for matrices & \texttt{(3,4) abs}\\
\texttt{sqrt} & \texttt{x -{}-{} sqrt(x)} & square root; a negative real promotes to complex & \texttt{-{}4 sqrt}\\
\texttt{\%} & \texttt{y x -{}-{} y y*x/100} & percent of y (HP \%; keeps y) & \texttt{200 25 \%}\\
\texttt{\%t} & \texttt{y x -{}-{} y 100*x/y} & x as a percent of the total y (HP \%T; keeps y) & \texttt{200 50 \%t}\\
\texttt{\%chg} & \texttt{y x -{}-{} y 100*(x-{}y)/y} & percent change from y to x (HP delta\%; keeps y) & \texttt{80 100 \%chg}\\
\texttt{mod} & \texttt{y x -{}-{} mod(y,x)} & floored modulo (result has the sign of x) & \texttt{-{}7 3 mod}\\
\texttt{fact} & \texttt{x -{}-{} x!} & factorial; Gamma(x+1) for non-{}integers & \texttt{5 fact}\\
\texttt{ncr} & \texttt{n k -{}-{} C(n,k)} & combinations & \texttt{52 5 ncr}\\
\texttt{npr} & \texttt{n k -{}-{} P(n,k)} & permutations & \texttt{10 3 npr}\\
\texttt{frac} & \texttt{x -{}-{} frac(x)} & fractional part, x -{} intg(x), with the sign of x & \texttt{3.7 frac}\\
\texttt{intg} & \texttt{x -{}-{} int(x)} & integer part, truncated toward zero as on HP calculators (-{}3.7 gives -{}3) & \texttt{3.7 intg}\\
\end{longtable}

\subsubsection*{Exp/Log}
\begin{longtable}{@{}>{\raggedright\arraybackslash}p{0.12\textwidth}>{\raggedright\arraybackslash}p{0.19\textwidth}>{\raggedright\arraybackslash}p{0.36\textwidth}>{\raggedright\arraybackslash}p{0.23\textwidth}@{}}
\toprule \textbf{Command} & \textbf{Stack} & \textbf{Description} & \textbf{Example}\\ \midrule \endfirsthead
\toprule \textbf{Command} & \textbf{Stack} & \textbf{Description} & \textbf{Example}\\ \midrule \endhead
\bottomrule \endlastfoot
\texttt{exp} & \texttt{x -{}-{} e\^{}x} & exponential (elementwise for matrices) & \texttt{1 exp}\\
\texttt{ln} & \texttt{x -{}-{} ln(x)} & natural logarithm; a negative real promotes to complex & \texttt{10 ln}\\
\texttt{log} & \texttt{x -{}-{} log10(x)} & base-{}10 logarithm; a negative real promotes to complex & \texttt{100 log}\\
\end{longtable}

\subsubsection*{Trigonometry}
\begin{longtable}{@{}>{\raggedright\arraybackslash}p{0.12\textwidth}>{\raggedright\arraybackslash}p{0.19\textwidth}>{\raggedright\arraybackslash}p{0.36\textwidth}>{\raggedright\arraybackslash}p{0.23\textwidth}@{}}
\toprule \textbf{Command} & \textbf{Stack} & \textbf{Description} & \textbf{Example}\\ \midrule \endfirsthead
\toprule \textbf{Command} & \textbf{Stack} & \textbf{Description} & \textbf{Example}\\ \midrule \endhead
\bottomrule \endlastfoot
\texttt{sin} & \texttt{z -{}-{} sin(z)} & sine of a real or complex argument (radians); elementwise for matrices & \texttt{0.5 sin}\\
\texttt{cos} & \texttt{z -{}-{} cos(z)} & cosine of real or complex argument & \texttt{1 cos}\\
\texttt{tan} & \texttt{z -{}-{} tan(z)} & tangent of real or complex argument & \texttt{0.3 tan}\\
\texttt{asin} & \texttt{z -{}-{} asin(z)} & inverse sine (principal value) & \texttt{0.5 asin}\\
\texttt{acos} & \texttt{z -{}-{} acos(z)} & inverse cosine (principal value) & \texttt{0.5 acos}\\
\texttt{atan} & \texttt{z -{}-{} atan(z)} & inverse tangent (principal value) & \texttt{1 atan}\\
\texttt{sinh} & \texttt{z -{}-{} sinh(z)} & hyperbolic sine & \texttt{1 sinh}\\
\texttt{cosh} & \texttt{z -{}-{} cosh(z)} & hyperbolic cosine & \texttt{1 cosh}\\
\texttt{tanh} & \texttt{z -{}-{} tanh(z)} & hyperbolic tangent & \texttt{1 tanh}\\
\texttt{asinh} & \texttt{z -{}-{} asinh(z)} & inverse hyperbolic sine & \texttt{1 asinh}\\
\texttt{acosh} & \texttt{z -{}-{} acosh(z)} & inverse hyperbolic cosine & \texttt{2 acosh}\\
\texttt{atanh} & \texttt{z -{}-{} atanh(z)} & inverse hyperbolic tangent & \texttt{0.5 atanh}\\
\end{longtable}

\subsubsection*{Constants}
\begin{longtable}{@{}>{\raggedright\arraybackslash}p{0.12\textwidth}>{\raggedright\arraybackslash}p{0.19\textwidth}>{\raggedright\arraybackslash}p{0.36\textwidth}>{\raggedright\arraybackslash}p{0.23\textwidth}@{}}
\toprule \textbf{Command} & \textbf{Stack} & \textbf{Description} & \textbf{Example}\\ \midrule \endfirsthead
\toprule \textbf{Command} & \textbf{Stack} & \textbf{Description} & \textbf{Example}\\ \midrule \endhead
\bottomrule \endlastfoot
\texttt{pi} & \texttt{-{}-{} π} & push π & \texttt{pi}\\
\texttt{e} & \texttt{-{}-{} e} & push Euler\textquotesingle{}s number & \texttt{e}\\
\texttt{gravity} & \texttt{-{}-{} g} & standard gravity, 9.80665 m/s\^{}2 & \texttt{gravity}\\
\texttt{inf} & \texttt{-{}-{} +∞} & positive infinity & \texttt{inf}\\
\texttt{nan} & \texttt{-{}-{} NaN} & not-{}a-{}number & \texttt{nan}\\
\end{longtable}

\subsubsection*{Logic}
\begin{longtable}{@{}>{\raggedright\arraybackslash}p{0.12\textwidth}>{\raggedright\arraybackslash}p{0.19\textwidth}>{\raggedright\arraybackslash}p{0.36\textwidth}>{\raggedright\arraybackslash}p{0.23\textwidth}@{}}
\toprule \textbf{Command} & \textbf{Stack} & \textbf{Description} & \textbf{Example}\\ \midrule \endfirsthead
\toprule \textbf{Command} & \textbf{Stack} & \textbf{Description} & \textbf{Example}\\ \midrule \endhead
\bottomrule \endlastfoot
\texttt{eq} & \texttt{a b -{}-{} flag} & 1 if a equals b, else 0; numbers, strings, or elementwise for matrices & \texttt{2 2 eq}\\
\texttt{neq} & \texttt{a b -{}-{} flag} & 1 if a differs from b; numbers, strings, or elementwise for matrices & \texttt{2 3 neq}\\
\texttt{lt} & \texttt{a b -{}-{} flag} & 1 if a \textless{} b (complex numbers compare by magnitude); elementwise for matrices & \texttt{2 3 lt}\\
\texttt{leq} & \texttt{a b -{}-{} flag} & 1 if a \textless{}= b; elementwise for matrices & \texttt{2 3 leq}\\
\texttt{gt} & \texttt{a b -{}-{} flag} & 1 if a \textgreater{} b; elementwise for matrices & \texttt{3 2 gt}\\
\texttt{geq} & \texttt{a b -{}-{} flag} & 1 if a \textgreater{}= b; elementwise for matrices & \texttt{3 2 geq}\\
\texttt{and} & \texttt{a b -{}-{} flag} & logical and (0 is false, anything else true) & \texttt{1 0 and}\\
\texttt{or} & \texttt{a b -{}-{} flag} & logical or & \texttt{1 0 or}\\
\texttt{not} & \texttt{a -{}-{} flag} & logical not & \texttt{0 not}\\
\end{longtable}

\subsubsection*{Stack}
\begin{longtable}{@{}>{\raggedright\arraybackslash}p{0.12\textwidth}>{\raggedright\arraybackslash}p{0.19\textwidth}>{\raggedright\arraybackslash}p{0.36\textwidth}>{\raggedright\arraybackslash}p{0.23\textwidth}@{}}
\toprule \textbf{Command} & \textbf{Stack} & \textbf{Description} & \textbf{Example}\\ \midrule \endfirsthead
\toprule \textbf{Command} & \textbf{Stack} & \textbf{Description} & \textbf{Example}\\ \midrule \endhead
\bottomrule \endlastfoot
\texttt{dup} & \texttt{x -{}-{} x x} & duplicate top of stack & \texttt{5 dup}\\
\texttt{drop} & \texttt{x -{}-{}} & drop the top stack element & \texttt{1 2 drop}\\
\texttt{swap} & \texttt{a b -{}-{} b a} & swap top two elements & \texttt{1 2 swap}\\
\texttt{over} & \texttt{a b -{}-{} a b a} & copy second item to top & \texttt{1 2 over}\\
\texttt{nip} & \texttt{a b -{}-{} b} & drop second-{}from-{}top element (keep top) & \texttt{1 2 nip}\\
\texttt{tuck} & \texttt{a b -{}-{} b a b} & duplicate second item and tuck below top & \texttt{1 2 tuck}\\
\texttt{roll} & \texttt{... n -{}-{} ...} & move the item at depth n to the top (1 roll = swap, 2 roll = rot) & \texttt{1 2 3 2 roll}\\
\texttt{pick} & \texttt{... n -{}-{} ... v} & copy the item at depth n to the top (0 pick = dup) & \texttt{1 2 3 2 pick}\\
\texttt{depth} & \texttt{-{}-{} n} & number of items on the stack & \texttt{depth}\\
\texttt{clst} & \texttt{... -{}-{}} & clear the stack & \texttt{clst}\\
\texttt{savestack} & \texttt{-{}-{}} & save the whole stack to stack\_backup.dat in the config directory & \texttt{savestack}\\
\texttt{loadstack} & \texttt{-{}-{} ...} & push the stack saved by savestack & \texttt{loadstack}\\
\end{longtable}

\subsubsection*{Display}
\begin{longtable}{@{}>{\raggedright\arraybackslash}p{0.12\textwidth}>{\raggedright\arraybackslash}p{0.19\textwidth}>{\raggedright\arraybackslash}p{0.36\textwidth}>{\raggedright\arraybackslash}p{0.23\textwidth}@{}}
\toprule \textbf{Command} & \textbf{Stack} & \textbf{Description} & \textbf{Example}\\ \midrule \endfirsthead
\toprule \textbf{Command} & \textbf{Stack} & \textbf{Description} & \textbf{Example}\\ \midrule \endhead
\bottomrule \endlastfoot
\texttt{print} & \texttt{x -{}-{} x} & print the top of the stack (a matrix in full) & \texttt{42 print}\\
\texttt{pm} & \texttt{A -{}-{} A} & print the matrix on top of the stack in full & \texttt{[2 2 \$ 1 2 3 4] pm}\\
\texttt{ps} & \texttt{-{}-{}} & print the whole stack & \texttt{ps}\\
\texttt{setprec} & \texttt{n -{}-{}} & digits after the decimal point (fixed) or significant digits (scientific), 0-{}20 & \texttt{10 setprec}\\
\texttt{sfs} & \texttt{-{}-{}} & toggle between fixed-{}point and scientific display & \texttt{sfs}\\
\end{longtable}

\subsubsection*{Words}
\begin{longtable}{@{}>{\raggedright\arraybackslash}p{0.12\textwidth}>{\raggedright\arraybackslash}p{0.19\textwidth}>{\raggedright\arraybackslash}p{0.36\textwidth}>{\raggedright\arraybackslash}p{0.23\textwidth}@{}}
\toprule \textbf{Command} & \textbf{Stack} & \textbf{Description} & \textbf{Example}\\ \midrule \endfirsthead
\toprule \textbf{Command} & \textbf{Stack} & \textbf{Description} & \textbf{Example}\\ \midrule \endhead
\bottomrule \endlastfoot
\texttt{listwords} & \texttt{-{}-{}} & list the user-{}defined words & \texttt{listwords}\\
\texttt{listmacros} & \texttt{-{}-{}} & list the words of the loaded word modules & \texttt{listmacros}\\
\texttt{loadwords} & \texttt{-{}-{}} & load the words saved in user\_words.txt (also done at startup) & \texttt{loadwords}\\
\texttt{savewords} & \texttt{-{}-{}} & save all user-{}defined words to user\_words.txt in the config directory & \texttt{savewords}\\
\texttt{delword} & \texttt{\textquotedbl{}name\textquotedbl{} -{}-{}} & delete a user-{}defined word & \texttt{\textquotedbl{}sq\textquotedbl{} delword}\\
\texttt{clrwords} & \texttt{-{}-{}} & delete all user-{}defined words & \texttt{clrwords}\\
\texttt{use} & \texttt{\textquotedbl{}name\textquotedbl{} -{}-{}} & load a word module: modules/name.mm in the config directory, or a path & \texttt{\textquotedbl{}finance\textquotedbl{} use}\\
\texttt{unuse} & \texttt{\textquotedbl{}name\textquotedbl{} -{}-{}} & unload a word module and its words & \texttt{\textquotedbl{}finance\textquotedbl{} unuse}\\
\end{longtable}

\subsubsection*{Session}
\begin{longtable}{@{}>{\raggedright\arraybackslash}p{0.12\textwidth}>{\raggedright\arraybackslash}p{0.19\textwidth}>{\raggedright\arraybackslash}p{0.36\textwidth}>{\raggedright\arraybackslash}p{0.23\textwidth}@{}}
\toprule \textbf{Command} & \textbf{Stack} & \textbf{Description} & \textbf{Example}\\ \midrule \endfirsthead
\toprule \textbf{Command} & \textbf{Stack} & \textbf{Description} & \textbf{Example}\\ \midrule \endhead
\bottomrule \endlastfoot
\texttt{load} & \texttt{\textquotedbl{}file.so\textquotedbl{} -{}-{}} & load a compiled module (a shared object exporting mm41\_module); stays loaded & \texttt{\textquotedbl{}hello.so\textquotedbl{} load}\\
\texttt{eval} & \texttt{s -{}-{}} & evaluate the string s as a line of input & \texttt{\textquotedbl{}1 2 +\textquotedbl{} eval}\\
\texttt{batch} & \texttt{s -{}-{}} & run the lines of file s as input & \texttt{\textquotedbl{}calc.txt\textquotedbl{} batch}\\
\texttt{help} & \texttt{-{}-{}} & print the quick-{}start help screen & \texttt{help}\\
\texttt{usage} & \texttt{\textquotedbl{}name\textquotedbl{} -{}-{}} & show the stack picture, description and example of a command, or a word\textquotesingle{}s definition & \texttt{\textquotedbl{}sin\textquotedbl{} usage}\\
\texttt{listfcns} & \texttt{-{}-{}} & list all commands by category & \texttt{listfcns}\\
\texttt{clrhist} & \texttt{-{}-{}} & clear the readline history & \texttt{clrhist}\\
\texttt{fuck} & \texttt{-{}-{}} & an Easter egg & \texttt{fuck}\\
\texttt{modules} & \texttt{-{}-{}} & list the loaded modules with their versions & \texttt{modules}\\
\end{longtable}

\subsection*{Module \texttt{decimal} 1.0 --- decimal numbers for money: number mode, conversions, rounding}
\subsubsection*{Decimal}
\begin{longtable}{@{}>{\raggedright\arraybackslash}p{0.12\textwidth}>{\raggedright\arraybackslash}p{0.19\textwidth}>{\raggedright\arraybackslash}p{0.36\textwidth}>{\raggedright\arraybackslash}p{0.23\textwidth}@{}}
\toprule \textbf{Command} & \textbf{Stack} & \textbf{Description} & \textbf{Example}\\ \midrule \endfirsthead
\toprule \textbf{Command} & \textbf{Stack} & \textbf{Description} & \textbf{Example}\\ \midrule \endhead
\bottomrule \endlastfoot
\texttt{decimal} & \texttt{-{}-{}} & number mode: literals are decimals (34 digits, half-{}even) until `binary` & \texttt{decimal 0.1 0.2 +}\\
\texttt{binary} & \texttt{-{}-{}} & number mode: literals are binary doubles (the default) & \texttt{binary}\\
\texttt{\textgreater{}dec} & \texttt{x -{}-{} d} & convert a real to a decimal by its shortest round-{}trip form (0.1 stays 0.1) & \texttt{0.1 \textgreater{}dec}\\
\texttt{\textgreater{}bin} & \texttt{d -{}-{} x} & convert a decimal to a binary real (needed before matrices and the numerical methods) & \texttt{1d 3d / \textgreater{}bin}\\
\texttt{round} & \texttt{x n -{}-{} d} & round to n decimal places, half-{}even; the result is a decimal & \texttt{2.345d 2 round}\\
\texttt{decmode} & \texttt{-{}-{} flag} & 1 if the number mode is decimal & \texttt{decmode}\\
\texttt{isdec} & \texttt{x -{}-{} flag} & 1 if x is a decimal & \texttt{1d isdec}\\
\end{longtable}

\subsection*{Module \texttt{finance} 1.0 --- time value of money, cash flows, depreciation, bonds, day counts}
\subsubsection*{Finance}
\begin{longtable}{@{}>{\raggedright\arraybackslash}p{0.12\textwidth}>{\raggedright\arraybackslash}p{0.19\textwidth}>{\raggedright\arraybackslash}p{0.36\textwidth}>{\raggedright\arraybackslash}p{0.23\textwidth}@{}}
\toprule \textbf{Command} & \textbf{Stack} & \textbf{Description} & \textbf{Example}\\ \midrule \endfirsthead
\toprule \textbf{Command} & \textbf{Stack} & \textbf{Description} & \textbf{Example}\\ \midrule \endhead
\bottomrule \endlastfoot
\texttt{n} & \texttt{x -{}-{}} & store the number of periods in the TVM register n & \texttt{12 n}\\
\texttt{i\%} & \texttt{x -{}-{}} & store the interest rate per period, in percent, in i\% & \texttt{1 i\%}\\
\texttt{pv} & \texttt{x -{}-{}} & store the present value (cash received +, paid -{}) & \texttt{10000 pv}\\
\texttt{pmt} & \texttt{x -{}-{}} & store the periodic payment & \texttt{-{}888.49 pmt}\\
\texttt{fv} & \texttt{x -{}-{}} & store the future value & \texttt{0 fv}\\
\texttt{beg} & \texttt{-{}-{}} & payments fall at the beginning of each period (annuity due) & \texttt{beg}\\
\texttt{end} & \texttt{-{}-{}} & payments fall at the end of each period (the default) & \texttt{end}\\
\texttt{clrtvm} & \texttt{-{}-{}} & clear the five TVM registers and select `end` & \texttt{clrtvm}\\
\texttt{tvm} & \texttt{-{}-{}} & show the TVM registers & \texttt{tvm}\\
\texttt{?n} & \texttt{-{}-{} n} & solve for the number of periods; the result is stored and pushed (a decimal) & \texttt{clrtvm 1 i\% 10000 pv -{}888.49 pmt 0 fv ?n}\\
\texttt{?i\%} & \texttt{-{}-{} i} & solve for the rate per period in percent (bracketing scan, then bisection) & \texttt{clrtvm 12 n 10000 pv -{}888.49 pmt 0 fv ?i\%}\\
\texttt{?pv} & \texttt{-{}-{} pv} & solve for the present value & \texttt{clrtvm 12 n 1 i\% -{}888.49 pmt 0 fv ?pv}\\
\texttt{?pmt} & \texttt{-{}-{} pmt} & solve for the periodic payment & \texttt{clrtvm 12 n 1 i\% 10000 pv 0 fv ?pmt}\\
\texttt{?fv} & \texttt{-{}-{} fv} & solve for the future value & \texttt{clrtvm 12 n 1 i\% 0 pv -{}100 pmt ?fv}\\
\texttt{amort} & \texttt{-{}-{}} & print the amortisation schedule of the TVM registers (n periods of pmt on pv at i\%) & \texttt{clrtvm 12 n 1 i\% 10000 pv 0 fv ?pmt drop amort}\\
\texttt{amortm} & \texttt{-{}-{} M} & the amortisation schedule as an n×4 real matrix: payment, interest, principal, balance & \texttt{clrtvm 12 n 1 i\% 10000 pv 0 fv ?pmt drop amortm}\\
\texttt{npv} & \texttt{cf rate -{}-{} x} & net present value of cash flows (a vector in time order, first at t=0, or N×2 amount/count) at rate percent per period & \texttt{[6 1 \$ -{}1000 300 300 300 300 300] 10 npv}\\
\texttt{irr} & \texttt{cf -{}-{} x} & internal rate of return in percent per period of a cash-{}flow vector & \texttt{[2 2 \$ -{}1000 1 300 5] irr}\\
\texttt{nom\textgreater{}eff} & \texttt{nom m -{}-{} eff} & nominal annual rate compounded m times a year to the effective annual rate (percent) & \texttt{12 12 nom\textgreater{}eff}\\
\texttt{eff\textgreater{}nom} & \texttt{eff m -{}-{} nom} & effective annual rate to the nominal rate compounded m times a year & \texttt{12.6825 12 eff\textgreater{}nom}\\
\texttt{sl} & \texttt{cost salvage life -{}-{} d} & straight-{}line depreciation per period & \texttt{10000 1000 5 sl}\\
\texttt{soyd} & \texttt{cost salvage life year -{}-{} d} & sum-{}of-{}the-{}years\textquotesingle{}-{}digits depreciation in the given year & \texttt{10000 1000 5 1 soyd}\\
\texttt{db} & \texttt{cost salvage life year factor -{}-{} d} & declining-{}balance depreciation in the given year (factor 2 = double declining), never below salvage & \texttt{10000 1000 5 1 2 db}\\
\texttt{bprice} & \texttt{cpn yld n redeem -{}-{} price} & bond price per 100 face on a coupon date: coupon and yield in percent per period, n periods, redemption value & \texttt{3 3.5 20 100 bprice}\\
\texttt{byield} & \texttt{price cpn n redeem -{}-{} yld} & bond yield in percent per period from the price on a coupon date & \texttt{92.89 3 20 100 byield}\\
\texttt{xnpv} & \texttt{M rate -{}-{} x} & net present value of dated cash flows (N×2 matrix: amount, date YYYYMMDD) at an annual rate in percent, actual/365 from the first date & \texttt{[5 2 \$ -{}10000 20080101 2750 20080301 4250 20081030 3250 20090215 2750 20090401] 9 xnpv}\\
\texttt{xirr} & \texttt{M -{}-{} x} & annual internal rate of return in percent of dated cash flows & \texttt{[3 2 \$ -{}1000 20250101 600 20250701 600 20260101] xirr}\\
\texttt{xprice} & \texttt{settle maturity cpn yld redeem freq basis -{}-{} price} & clean price per 100 face of a dated bond: annual coupon and yield in percent, coupons per year, day-{}count basis 0..4 & \texttt{20080215 20171115 5.75 6.5 100 2 0 xprice}\\
\texttt{xyield} & \texttt{settle maturity cpn price redeem freq basis -{}-{} yld} & annual yield in percent of a dated bond from its clean price & \texttt{20080215 20161115 5.75 95.04287 100 2 0 xyield}\\
\texttt{xaccr} & \texttt{settle maturity cpn freq basis -{}-{} accrued} & accrued interest per 100 face at settlement & \texttt{20080215 20171115 5.75 2 0 xaccr}\\
\texttt{yearfrac} & \texttt{d1 d2 basis -{}-{} f} & year fraction between two dates; basis 0 = 30/360 US, 1 = actual/actual, 2 = actual/360, 3 = actual/365, 4 = 30E/360 & \texttt{\textquotedbl{}1.1.2007\textquotedbl{} \textquotedbl{}1.1.2008\textquotedbl{} 3 yearfrac}\\
\texttt{date\textgreater{}ymd} & \texttt{date -{}-{} yyyymmdd} & a date string d.m.yyyy as the integer YYYYMMDD (for cash-{}flow matrices) & \texttt{\textquotedbl{}15.3.2026\textquotedbl{} date\textgreater{}ymd}\\
\texttt{ymd\textgreater{}date} & \texttt{yyyymmdd -{}-{} date} & an integer YYYYMMDD as the date string d.m.yyyy of the dates module & \texttt{20260315 ymd\textgreater{}date}\\
\end{longtable}

\subsection*{Module \texttt{complex} 1.0 --- complex numbers: parts, argument, conjugate, construction}
\subsubsection*{Complex}
\begin{longtable}{@{}>{\raggedright\arraybackslash}p{0.12\textwidth}>{\raggedright\arraybackslash}p{0.19\textwidth}>{\raggedright\arraybackslash}p{0.36\textwidth}>{\raggedright\arraybackslash}p{0.23\textwidth}@{}}
\toprule \textbf{Command} & \textbf{Stack} & \textbf{Description} & \textbf{Example}\\ \midrule \endfirsthead
\toprule \textbf{Command} & \textbf{Stack} & \textbf{Description} & \textbf{Example}\\ \midrule \endhead
\bottomrule \endlastfoot
\texttt{re} & \texttt{z -{}-{} Re(z)} & real part (a real is unchanged) & \texttt{(1,2) re}\\
\texttt{im} & \texttt{z -{}-{} Im(z)} & imaginary part (0 for a real) & \texttt{(1,2) im}\\
\texttt{arg} & \texttt{z -{}-{} arg(z)} & argument (phase) in radians, in (-{}pi, pi] & \texttt{(0,1) arg}\\
\texttt{conj} & \texttt{z -{}-{} conj(z)} & complex conjugate (elementwise for matrices; identity for reals) & \texttt{(1,2) conj}\\
\texttt{re2c} & \texttt{x -{}-{} (x,0)} & promote a real to complex & \texttt{5 re2c}\\
\texttt{split\_c} & \texttt{z -{}-{} Re(z) Im(z)} & split a complex number into real and imaginary parts & \texttt{(1,2) split\_c}\\
\texttt{j2r} & \texttt{re im -{}-{} (re,im)} & join two reals into a complex number (imaginary part on top) & \texttt{1 2 j2r}\\
\end{longtable}

\subsection*{Module \texttt{special} 1.0 --- gamma, beta and the standard normal distribution}
\subsubsection*{Special functions}
\begin{longtable}{@{}>{\raggedright\arraybackslash}p{0.12\textwidth}>{\raggedright\arraybackslash}p{0.19\textwidth}>{\raggedright\arraybackslash}p{0.36\textwidth}>{\raggedright\arraybackslash}p{0.23\textwidth}@{}}
\toprule \textbf{Command} & \textbf{Stack} & \textbf{Description} & \textbf{Example}\\ \midrule \endfirsthead
\toprule \textbf{Command} & \textbf{Stack} & \textbf{Description} & \textbf{Example}\\ \midrule \endhead
\bottomrule \endlastfoot
\texttt{npdf} & \texttt{x -{}-{} phi(x)} & standard normal density & \texttt{0 npdf}\\
\texttt{ncdf} & \texttt{x -{}-{} Phi(x)} & standard normal cumulative distribution & \texttt{1.96 ncdf}\\
\texttt{nquant} & \texttt{p -{}-{} x} & standard normal quantile (inverse of ncdf), 0 \textless{} p \textless{} 1 & \texttt{0.975 nquant}\\
\texttt{gamma} & \texttt{x -{}-{} Gamma(x)} & gamma function & \texttt{5 gamma}\\
\texttt{ln\_gamma} & \texttt{x -{}-{} lnΓ(x)} & log gamma function & \texttt{10 ln\_gamma}\\
\texttt{beta} & \texttt{a b -{}-{} B(a,b)} & beta function & \texttt{2 3 beta}\\
\texttt{ln\_beta} & \texttt{a b -{}-{} lnB(a,b)} & log beta function & \texttt{2 3 ln\_beta}\\
\end{longtable}

\subsection*{Module \texttt{matrices} 1.0 --- linear algebra, matrix construction and reduction, polynomials}
\subsubsection*{Matrices}
\begin{longtable}{@{}>{\raggedright\arraybackslash}p{0.12\textwidth}>{\raggedright\arraybackslash}p{0.19\textwidth}>{\raggedright\arraybackslash}p{0.36\textwidth}>{\raggedright\arraybackslash}p{0.23\textwidth}@{}}
\toprule \textbf{Command} & \textbf{Stack} & \textbf{Description} & \textbf{Example}\\ \midrule \endfirsthead
\toprule \textbf{Command} & \textbf{Stack} & \textbf{Description} & \textbf{Example}\\ \midrule \endhead
\bottomrule \endlastfoot
\texttt{tran} & \texttt{A -{}-{} A\^{}T} & transpose & \texttt{[2 3 \$ 1 2 3 4 5 6] tran}\\
\texttt{\textquotesingle{}} & \texttt{A -{}-{} A\textquotesingle{}} & transpose (alias of tran) & \texttt{[2 3 \$ 1 2 3 4 5 6] \textquotesingle{}}\\
\texttt{minv} & \texttt{A -{}-{} A\^{}-{}1} & inverse of a square real or complex matrix (error if singular) & \texttt{[2 2 \$ 4 7 2 6] minv}\\
\texttt{pinv} & \texttt{A -{}-{} A\^{}+} & Moore-{}Penrose pseudoinverse of a square real matrix (SVD) & \texttt{[2 2 \$ 1 2 2 4] pinv}\\
\texttt{det} & \texttt{A -{}-{} det(A)} & determinant of a square real or complex matrix & \texttt{[2 2 \$ 4 7 2 6] det}\\
\texttt{eig} & \texttt{A -{}-{} V L} & eigenvectors (columns of V) and eigenvalues (diagonal of L) of a square real matrix; both complex & \texttt{[2 2 \$ 2 1 1 2] eig}\\
\texttt{chol} & \texttt{A -{}-{} L} & Cholesky factor of a symmetric positive definite real matrix, A = L L\^{}T & \texttt{[2 2 \$ 4 2 2 3] chol}\\
\texttt{svd} & \texttt{A -{}-{} U S V} & singular value decomposition of a real matrix, A = U S V\^{}T & \texttt{[2 2 \$ 4 7 2 6] svd}\\
\texttt{kron} & \texttt{A B -{}-{} A(x)B} & Kronecker product & \texttt{[2 2 \$ 1 2 3 4] [1 2 \$ 5 6] kron}\\
\texttt{dim} & \texttt{A -{}-{} A rows cols} & push the dimensions (the matrix stays) & \texttt{[2 3 \$ 1 2 3 4 5 6] dim}\\
\texttt{reshape} & \texttt{A rows cols -{}-{} B} & reshape in row-{}major order; rows*cols must match & \texttt{[2 3 \$ 1 2 3 4 5 6] 3 2 reshape}\\
\texttt{get\_aij} & \texttt{A i j -{}-{} A a\_ij} & push element (i,j), 0-{}based (the matrix stays) & \texttt{[2 2 \$ 1 2 3 4] 0 1 get\_aij}\\
\texttt{set\_aij} & \texttt{A i j x -{}-{} A} & set element (i,j), 0-{}based, in place; a complex x promotes a real matrix & \texttt{[2 2 \$ 1 2 3 4] 0 1 9 set\_aij}\\
\texttt{split\_mat} & \texttt{A -{}-{} a11 a12 ...} & push all elements as scalars, row by row & \texttt{[2 2 \$ 1 2 3 4] split\_mat}\\
\texttt{diag} & \texttt{A -{}-{} v} & main diagonal as a row vector & \texttt{[2 2 \$ 1 2 3 4] diag}\\
\texttt{to\_diag} & \texttt{v -{}-{} D} & diagonal matrix from a row or column vector & \texttt{[1 3 \$ 1 2 3] to\_diag}\\
\texttt{eye} & \texttt{n -{}-{} I} & n x n identity matrix & \texttt{3 eye}\\
\texttt{ones} & \texttt{rows cols -{}-{} A} & matrix filled with ones & \texttt{2 3 ones}\\
\texttt{zeroes} & \texttt{rows cols -{}-{} A} & matrix filled with zeros & \texttt{2 3 zeroes}\\
\texttt{rand} & \texttt{rows cols -{}-{} A} & matrix of uniform(0,1) randoms & \texttt{2 2 rand}\\
\texttt{randn} & \texttt{rows cols -{}-{} A} & matrix of standard normal randoms & \texttt{2 2 randn}\\
\texttt{rrange} & \texttt{start step end -{}-{} v} & row vector start, start+step, ... up to end & \texttt{0 0.25 1 rrange}\\
\texttt{join\_v} & \texttt{A B -{}-{} [A;B]} & stack two matrices vertically (same number of columns) & \texttt{[1 2 \$ 1 2] [1 2 \$ 3 4] join\_v}\\
\texttt{join\_h} & \texttt{A B -{}-{} [A\textbar{}B]} & join two matrices side by side (same number of rows) & \texttt{[2 1 \$ 1 2] [2 1 \$ 3 4] join\_h}\\
\texttt{cumsum\_r} & \texttt{A -{}-{} B} & cumulative sum along each row & \texttt{[2 2 \$ 1 2 3 4] cumsum\_r}\\
\texttt{cumsum\_c} & \texttt{A -{}-{} B} & cumulative sum down each column & \texttt{[2 2 \$ 1 2 3 4] cumsum\_c}\\
\texttt{cmean} & \texttt{A -{}-{} v} & mean of each column (row vector) & \texttt{[2 2 \$ 1 2 3 4] cmean}\\
\texttt{rmean} & \texttt{A -{}-{} v} & mean of each row (column vector) & \texttt{[2 2 \$ 1 2 3 4] rmean}\\
\texttt{csum} & \texttt{A -{}-{} v} & sum of each column (row vector) & \texttt{[2 2 \$ 1 2 3 4] csum}\\
\texttt{rsum} & \texttt{A -{}-{} v} & sum of each row (column vector) & \texttt{[2 2 \$ 1 2 3 4] rsum}\\
\texttt{cvar} & \texttt{A -{}-{} v} & sample variance of each column (row vector) & \texttt{[2 2 \$ 1 2 3 4] cvar}\\
\texttt{rvar} & \texttt{A -{}-{} v} & sample variance of each row (column vector) & \texttt{[2 2 \$ 1 2 3 4] rvar}\\
\texttt{cmin} & \texttt{A -{}-{} v} & minimum of each column & \texttt{[2 2 \$ 1 2 3 4] cmin}\\
\texttt{rmin} & \texttt{A -{}-{} v} & minimum of each row & \texttt{[2 2 \$ 1 2 3 4] rmin}\\
\texttt{cmax} & \texttt{A -{}-{} v} & maximum of each column & \texttt{[2 2 \$ 1 2 3 4] cmax}\\
\texttt{rmax} & \texttt{A -{}-{} v} & maximum of each row & \texttt{[2 2 \$ 1 2 3 4] rmax}\\
\end{longtable}

\subsubsection*{Polynomials}
\begin{longtable}{@{}>{\raggedright\arraybackslash}p{0.12\textwidth}>{\raggedright\arraybackslash}p{0.19\textwidth}>{\raggedright\arraybackslash}p{0.36\textwidth}>{\raggedright\arraybackslash}p{0.23\textwidth}@{}}
\toprule \textbf{Command} & \textbf{Stack} & \textbf{Description} & \textbf{Example}\\ \midrule \endfirsthead
\toprule \textbf{Command} & \textbf{Stack} & \textbf{Description} & \textbf{Example}\\ \midrule \endhead
\bottomrule \endlastfoot
\texttt{roots} & \texttt{coeffs -{}-{} roots} & roots of a real polynomial; coefficients highest degree first, result is a complex row vector & \texttt{[1 3 \$ 1 -{}3 2] roots}\\
\texttt{pval} & \texttt{coeffs x -{}-{} y} & evaluate a polynomial (coefficients highest degree first, real or complex) at x & \texttt{[1 3 \$ 1 0 -{}1] 2 pval}\\
\end{longtable}

\subsection*{Module \texttt{numerics} 1.0 --- integration, roots, minima, derivatives and sums over any word}
\subsubsection*{Numerics}
\begin{longtable}{@{}>{\raggedright\arraybackslash}p{0.12\textwidth}>{\raggedright\arraybackslash}p{0.19\textwidth}>{\raggedright\arraybackslash}p{0.36\textwidth}>{\raggedright\arraybackslash}p{0.23\textwidth}@{}}
\toprule \textbf{Command} & \textbf{Stack} & \textbf{Description} & \textbf{Example}\\ \midrule \endfirsthead
\toprule \textbf{Command} & \textbf{Stack} & \textbf{Description} & \textbf{Example}\\ \midrule \endhead
\bottomrule \endlastfoot
\texttt{exec} & \texttt{xt -{}-{}} & run an execution token (@name pushes one) & \texttt{3 @sq exec}\\
\texttt{integrate} & \texttt{a b xt -{}-{} I} & integral of f over [a,b] (tanh-{}sinh with adaptive bisection; handles endpoint singularities) & \texttt{0 pi @sin integrate}\\
\texttt{romberg} & \texttt{a b xt -{}-{} I} & integral of f over [a,b] by Romberg extrapolation & \texttt{0 1 @exp romberg}\\
\texttt{fzero} & \texttt{a b xt -{}-{} x} & root of f in the bracket [a,b] where f changes sign (Brent) & \texttt{2 4 @sin fzero}\\
\texttt{solve} & \texttt{guess xt -{}-{} x} & root of f near one guess (bracket search, then Brent) & \texttt{3 @sin solve}\\
\texttt{minimize} & \texttt{a b xt -{}-{} x} & argument of the minimum of f in [a,b] (Brent) & \texttt{0 5 @par minimize}\\
\texttt{deriv} & \texttt{x xt -{}-{} f\textquotesingle{}(x)} & derivative of f at x (Ridders extrapolation) & \texttt{2 @exp deriv}\\
\texttt{fsum} & \texttt{lo hi xt -{}-{} S} & sum of f(i) for integer i from lo to hi & \texttt{1 1000 @invsq fsum}\\
\texttt{fprod} & \texttt{lo hi xt -{}-{} P} & product of f(i) for integer i from lo to hi & \texttt{1 10 @ident fprod}\\
\texttt{tol} & \texttt{x -{}-{}} & set the convergence tolerance (default 1e-{}10) & \texttt{1e-{}6 tol}\\
\texttt{nevals} & \texttt{-{}-{} n} & function evaluations used by the last numerical call & \texttt{nevals}\\
\texttt{interr} & \texttt{-{}-{} e} & error estimate of the last integrate & \texttt{interr}\\
\end{longtable}

\subsection*{Module \texttt{strings} 1.0 --- string construction and manipulation}
\subsubsection*{Strings}
\begin{longtable}{@{}>{\raggedright\arraybackslash}p{0.12\textwidth}>{\raggedright\arraybackslash}p{0.19\textwidth}>{\raggedright\arraybackslash}p{0.36\textwidth}>{\raggedright\arraybackslash}p{0.23\textwidth}@{}}
\toprule \textbf{Command} & \textbf{Stack} & \textbf{Description} & \textbf{Example}\\ \midrule \endfirsthead
\toprule \textbf{Command} & \textbf{Stack} & \textbf{Description} & \textbf{Example}\\ \midrule \endhead
\bottomrule \endlastfoot
\texttt{scon} & \texttt{s1 s2 -{}-{} s1s2} & concatenate two strings & \texttt{\textquotedbl{}foo\textquotedbl{} \textquotedbl{}bar\textquotedbl{} scon}\\
\texttt{substr} & \texttt{s start end -{}-{} sub} & substring from index start (0-{}based) up to end (exclusive) & \texttt{\textquotedbl{}abcdef\textquotedbl{} 1 3 substr}\\
\texttt{s2l} & \texttt{s -{}-{} s\_lower} & convert string to lowercase & \texttt{\textquotedbl{}Hello\textquotedbl{} s2l}\\
\texttt{s2u} & \texttt{s -{}-{} s\_upper} & convert string to uppercase & \texttt{\textquotedbl{}Hello\textquotedbl{} s2u}\\
\texttt{slen} & \texttt{s -{}-{} n} & length of a string in bytes & \texttt{\textquotedbl{}hello\textquotedbl{} slen}\\
\texttt{srev} & \texttt{s -{}-{} s\_rev} & reverse string & \texttt{\textquotedbl{}abc\textquotedbl{} srev}\\
\texttt{int2str} & \texttt{x -{}-{} s} & integer part of x as a string & \texttt{123 int2str}\\
\end{longtable}

\subsection*{Module \texttt{registers} 1.0 --- 64 storage registers of any type}
\subsubsection*{Registers}
\begin{longtable}{@{}>{\raggedright\arraybackslash}p{0.12\textwidth}>{\raggedright\arraybackslash}p{0.19\textwidth}>{\raggedright\arraybackslash}p{0.36\textwidth}>{\raggedright\arraybackslash}p{0.23\textwidth}@{}}
\toprule \textbf{Command} & \textbf{Stack} & \textbf{Description} & \textbf{Example}\\ \midrule \endfirsthead
\toprule \textbf{Command} & \textbf{Stack} & \textbf{Description} & \textbf{Example}\\ \midrule \endhead
\bottomrule \endlastfoot
\texttt{sto} & \texttt{x n -{}-{}} & store x in register n (0-{}63); any type & \texttt{42 0 sto}\\
\texttt{rcl} & \texttt{n -{}-{} x} & push a copy of register n & \texttt{0 rcl}\\
\texttt{ffr} & \texttt{-{}-{} n} & index of the first free register & \texttt{ffr}\\
\texttt{pr} & \texttt{-{}-{}} & show which registers are occupied & \texttt{pr}\\
\texttt{saveregs} & \texttt{-{}-{}} & save all registers to registers.txt in the config directory & \texttt{saveregs}\\
\texttt{loadregs} & \texttt{-{}-{}} & load the registers saved by saveregs & \texttt{loadregs}\\
\texttt{clregs} & \texttt{-{}-{}} & clear all registers & \texttt{clregs}\\
\end{longtable}

\subsection*{Module \texttt{dates} 1.1 --- calendar arithmetic, workdays and holidays, time zones, sunrise and sunset}
\subsubsection*{Dates}
\begin{longtable}{@{}>{\raggedright\arraybackslash}p{0.12\textwidth}>{\raggedright\arraybackslash}p{0.19\textwidth}>{\raggedright\arraybackslash}p{0.36\textwidth}>{\raggedright\arraybackslash}p{0.23\textwidth}@{}}
\toprule \textbf{Command} & \textbf{Stack} & \textbf{Description} & \textbf{Example}\\ \midrule \endfirsthead
\toprule \textbf{Command} & \textbf{Stack} & \textbf{Description} & \textbf{Example}\\ \midrule \endhead
\bottomrule \endlastfoot
\texttt{today} & \texttt{-{}-{} date} & today\textquotesingle{}s date as a string d.m.yyyy & \texttt{today}\\
\texttt{now} & \texttt{-{}-{} \textquotedbl{}HH:MM\textquotedbl{}} & push the current local time as a string, same form as sunrise/sunset & \texttt{now}\\
\texttt{tz\_offset} & \texttt{date zone -{}-{} hours} & UTC offset in hours of an IANA zone on the given date, DST included & \texttt{\textquotedbl{}1.12.2026\textquotedbl{} \textquotedbl{}America/New\_York\textquotedbl{} tz\_offset}\\
\texttt{ddays} & \texttt{d1 d2 -{}-{} n} & days from date d1 to date d2 (d2 -{} d1) & \texttt{\textquotedbl{}1.1.2025\textquotedbl{} \textquotedbl{}31.12.2025\textquotedbl{} ddays}\\
\texttt{dateplus} & \texttt{date n -{}-{} date2} & the date n days later (n may be negative) & \texttt{\textquotedbl{}1.1.2025\textquotedbl{} 30 dateplus}\\
\texttt{dow} & \texttt{date -{}-{} name} & weekday name of a date & \texttt{\textquotedbl{}1.1.2025\textquotedbl{} dow}\\
\texttt{edmy} & \texttt{date -{}-{} y m d} & split a date into year, month, day (day on top) & \texttt{\textquotedbl{}1.1.2025\textquotedbl{} edmy}\\
\texttt{num2date} & \texttt{d m y -{}-{} date} & build a date string from day, month, year & \texttt{1 1 2025 num2date}\\
\texttt{days2eoy} & \texttt{date -{}-{} n} & days from the date to the end of its year & \texttt{\textquotedbl{}1.1.2025\textquotedbl{} days2eoy}\\
\texttt{sunrise} & \texttt{date lat lon zone -{}-{} time} & local sunrise HH:MM; zone is an IANA name or a UTC offset in hours & \texttt{\textquotedbl{}21.6.2025\textquotedbl{} 38.9 -{}77.0 \textquotedbl{}America/New\_York\textquotedbl{} sunrise}\\
\texttt{sunset} & \texttt{date lat lon zone -{}-{} time} & local sunset HH:MM & \texttt{\textquotedbl{}21.6.2025\textquotedbl{} 38.9 -{}77.0 -{}4 sunset}\\
\texttt{dawn} & \texttt{date lat lon zone -{}-{} time} & civil dawn HH:MM (sun 6 degrees below the horizon) & \texttt{today 38.9 -{}77.0 \textquotedbl{}America/New\_York\textquotedbl{} dawn}\\
\texttt{dusk} & \texttt{date lat lon zone -{}-{} time} & civil dusk HH:MM & \texttt{today 38.9 -{}77.0 \textquotedbl{}America/New\_York\textquotedbl{} dusk}\\
\end{longtable}

\subsubsection*{Workdays}
\begin{longtable}{@{}>{\raggedright\arraybackslash}p{0.12\textwidth}>{\raggedright\arraybackslash}p{0.19\textwidth}>{\raggedright\arraybackslash}p{0.36\textwidth}>{\raggedright\arraybackslash}p{0.23\textwidth}@{}}
\toprule \textbf{Command} & \textbf{Stack} & \textbf{Description} & \textbf{Example}\\ \midrule \endfirsthead
\toprule \textbf{Command} & \textbf{Stack} & \textbf{Description} & \textbf{Example}\\ \midrule \endhead
\bottomrule \endlastfoot
\texttt{workdays} & \texttt{d1 d2 -{}-{} n} & workdays from d1 to d2, both inclusive (Excel NETWORKDAYS): Mondays to Fridays not in the holiday calendar; negative if d2 \textless{} d1 & \texttt{\textquotedbl{}1.1.2026\textquotedbl{} \textquotedbl{}31.12.2026\textquotedbl{} workdays}\\
\texttt{workdayplus} & \texttt{date n -{}-{} date2} & the date n workdays after date (n may be negative; the start date is not counted, as Excel WORKDAY) & \texttt{\textquotedbl{}24.12.2026\textquotedbl{} 10 workdayplus}\\
\texttt{isworkday} & \texttt{date -{}-{} flag} & 1 if the date is a Monday to Friday and not a holiday & \texttt{20261126 isworkday}\\
\texttt{holidays} & \texttt{name -{}-{} H} & the holiday calendar as a vector of YYYYMMDD: the file holidays/name.txt in the config directory, or the built-{}in rules \textquotedbl{}us\textquotedbl{} (federal, observed) and \textquotedbl{}si\textquotedbl{} (Slovenia) for the years around today & \texttt{\textquotedbl{}us\textquotedbl{} holidays}\\
\texttt{yholidays} & \texttt{year name -{}-{} H} & the built-{}in rules for one year & \texttt{2027 \textquotedbl{}si\textquotedbl{} yholidays}\\
\texttt{sethol} & \texttt{H -{}-{}} & install a vector of dates YYYYMMDD as the holiday calendar for this session & \texttt{[2 1 \$ 20261126 20261225] sethol}\\
\texttt{usehol} & \texttt{name -{}-{}} & install the calendar `name` (file or rules) and remember it in config.txt & \texttt{\textquotedbl{}us\textquotedbl{} usehol}\\
\texttt{clrhol} & \texttt{-{}-{}} & remove the holiday calendar (weekends only) and forget it in config.txt & \texttt{clrhol}\\
\texttt{hol} & \texttt{-{}-{} H} & the installed holiday calendar as a vector & \texttt{hol}\\
\end{longtable}

\subsection*{Module \texttt{programs} 1.0 --- HP-41 style programs with labels, jumps and subroutines}
\subsubsection*{Programs}
\begin{longtable}{@{}>{\raggedright\arraybackslash}p{0.12\textwidth}>{\raggedright\arraybackslash}p{0.19\textwidth}>{\raggedright\arraybackslash}p{0.36\textwidth}>{\raggedright\arraybackslash}p{0.23\textwidth}@{}}
\toprule \textbf{Command} & \textbf{Stack} & \textbf{Description} & \textbf{Example}\\ \midrule \endfirsthead
\toprule \textbf{Command} & \textbf{Stack} & \textbf{Description} & \textbf{Example}\\ \midrule \endhead
\bottomrule \endlastfoot
\texttt{run} & \texttt{s -{}-{}} & load and run the HP-{}41 style program in file s (LBL/GOTO/GOSUB/RTN/END) & \texttt{\textquotedbl{}prog1.prg\textquotedbl{} run}\\
\end{longtable}

